
\documentclass[aps,preprint,showpacs,showkeys]{revtex4}%
\usepackage{amsfonts}
\usepackage{amsmath}
\usepackage[dvips]{graphicx}
\usepackage{amssymb}%}

\setcounter{MaxMatrixCols}{10}
%TCIDATA{OutputFilter=LATEX.DLL}
%TCIDATA{Version=5.50.0.2960}
%TCIDATA{<META NAME="SaveForMode" CONTENT="1">}
%TCIDATA{BibliographyScheme=Manual}
%TCIDATA{LastRevised=Friday, August 21, 2026 16:30:34}
%TCIDATA{<META NAME="GraphicsSave" CONTENT="32">}

\newtheorem{axiom}{Axiom}
\newtheorem{theorem}{Theorem}
\newtheorem{definition}{Definition}
\newtheorem{corollary}{Corollary}
\newtheorem{example}{Example}
\newtheorem{lemma}{Lemma}
\hypersetup{
    colorlinks=true,
    linkcolor=blue,
    citecolor=blue,
    urlcolor=blue,
}
\input{tcilatex}
\begin{document}

\title{The Principle of Reality: A Causal-Ontological Foundation for Physics}
\author{The Originator}
\date{August 2026}

\begin{abstract}
We present a formal framework for physical ontology based on the \textit{%
Principle of Reality}: an entity exists in physical reality iff it can be
connected to a measuring apparatus via a finite causal chain, with each
interaction consuming finite energy and finite time. This principle unifies
ontology, epistemology, and causality into a single consistent framework. We
show that quantum mechanics is not a theory of physical reality itself, but
a theory of \textit{information about} physical reality. The wavefunction,
Hilbert space, and operators are mathematical tools for encoding knowledge
about possible measurement outcomes.

We extend the formal analysis to include the fundamental wave-measurement
basis of all physical measurement: length measurements rely on wave
interference with energy $E \ge \hbar c/\lambda$ (where $\lambda$ is the
wavelength), while time measurements rely on oscillation frequency with
energy $E = hf$. This wave basis provides the physical foundation for the
practical energy bounds and naturally yields the Planck scale as the
practical limit of measurability. The gravitational redshift of atomic
transitions near black holes is examined as a rigorous test case,
demonstrating that the observed transition energy goes to zero at the event
horizon, requiring infinite time to measure.

The framework yields a dynamic ontology---new measurement procedures can be
discovered, expanding the set of known entities---but the criterion for
existence remains fixed and invariant. Key consequences include: a
Planck-scale bound on measurability, causal closure of physical reality, the
non-Boolean (orthomodular) lattice structure of measurement propositions,
and resolution of the measurement problem by distinguishing ontological
updates (none occur) from epistemic updates (information changes upon
measurement).
\end{abstract}

\maketitle

\affiliation{Independent Scholar}

\section{Introduction}

The relationship between \emph{what exists} and \emph{what can be known} has
been a central question in philosophy and physics for centuries. In the 20th
century, the development of quantum mechanics deepened this tension: the
wavefunction appears to describe a physical system, yet its interpretation
remains contentious. Does it describe reality itself, or merely our
knowledge of reality?

The \textbf{Principle of Reality} (PR) offers a decisive answer. It asserts
that physical existence is \emph{equivalent to causal measurability}. An
entity or event is part of physical reality iff there exists a finite,
practical procedure to register its effects. This is not merely an
epistemological claim about what we can know---it is an ontological claim
about what \emph{is}.

The PR is distinguished from related views by three key features:

\begin{enumerate}
\item \textbf{Ontological commitment:} Unmeasurable entities do not exist
(they are not merely "unknowable").

\item \textbf{Practicality constraint:} Measurements must consume finite
energy and finite time---ruling out "in principle but never in practice"
procedures.

\item \textbf{Causal chain requirement:} Indirect measurements are allowed,
but the chain of causal interactions must be finite.
\end{enumerate}

This paper formalizes the PR, shows its consistency with quantum mechanics
when interpreted as a theory of information, derives its consequences, and
addresses objections. We also provide a rigorous analysis of the
wave-measurement basis of all physical measurement and examine gravitational
redshift as a test case.

\section{The Wave-Measurement Foundation}

Before presenting the formal axioms, we establish the physical basis of all
measurement. This is essential because the PR's practicality constraint
derives directly from how measurements are actually performed.

\subsection{Length Measurement by Wave Interference}

All high-precision measurements of length ultimately rely on wave
interference. To resolve a distance $\Delta x$, one requires a probe wave
with wavelength $\lambda$ such that: 
\begin{equation}
\lambda \le \Delta x
\end{equation}

By the de Broglie relation, the energy of the probe is: 
\begin{equation}
E = \frac{hc}{\lambda} \ge \frac{hc}{\Delta x}
\end{equation}

Thus, measuring smaller distances requires higher energy. This is not an
assumption---it is a direct consequence of wave-particle duality.

\subsection{Time Measurement by Oscillation Frequency}

All time measurements rely on counting oscillations of a periodic system. To
resolve a time interval $\Delta t$, one requires a clock with frequency $f$
such that: 
\begin{equation}
f \ge \frac{1}{\Delta t}
\end{equation}

By the Planck-Einstein relation, the energy of the clock transition is: 
\begin{equation}
E = hf \ge \frac{h}{\Delta t}
\end{equation}

Thus, measuring shorter time intervals requires higher energy.

\subsection{The Universal Nature of Wave Measurement}

These two relationships reveal a fundamental truth: 
\begin{equation}
\frame{All physical measurements are wave measurements}
\end{equation}

This is because the only way to interact with a physical system is through
the exchange of quanta (photons, electrons, etc.), which are waves. The
energy dependence of measurement is therefore universal: 
\begin{equation}
E \ge \max\left(\frac{hc}{\Delta x}, \frac{h}{\Delta t}\right)
\end{equation}

This leads directly to the practical bounds on measurability.

\section{Core Definitions and Axioms}

\subsection{Physical Reality}

\begin{definition}[Physical Reality]
Let $\mathcal{R}$ be the set of all entities, properties, and events that
can affect a measuring apparatus through a finite causal chain. An entity $E$
is an element of $\mathcal{R}$ iff there exists a sequence: 
\begin{equation}
E \rightarrow I_1 \rightarrow I_2 \rightarrow \cdots \rightarrow I_n
\rightarrow D
\end{equation}
where:

\begin{itemize}
\item Each $I_k$ is a causal interaction

\item $D$ is a measuring apparatus with a binary (yes/no) output

\item Each $I_k$ consumes finite energy $\Delta E_k < \infty$

\item Each $I_k$ takes finite time $\Delta t_k < \infty$

\item The total energy $\sum_k \Delta E_k \le E_{\text{available}}$ (the
energy accessible in the observable universe)
\end{itemize}
\end{definition}

\subsection{Measurement Procedure}

\begin{definition}[Valid Measurement Procedure]
A valid measurement procedure for entity $E$ is a finite sequence of
physical interactions satisfying the conditions in Definition 1. The
procedure is \emph{practical} if each step is achievable with currently
available or theoretically possible technology (subject to physical laws).
\end{definition}

\subsection{The Measurement Bound (Wave-Derived)}

\begin{theorem}[Measurement Bound]
There exists a minimal length $L_0 > 0$ and minimal time $T_0 > 0$ such that
no proposition in $\mathcal{L}$ can assert a length $L < L_0$ or a time $T <
T_0$.
\end{theorem}

\begin{proof}
To measure a length $L$, one must use a probe with wavelength $\lambda \le L$, requiring energy $E \ge \hbar c/L$. To measure a time $T$, one requires energy $E \ge \hbar/T$. If $L$ or $T$ is too small, the required energy exceeds either the total energy of the observable universe or the threshold for black hole formation. Therefore, no valid measurement procedure exists for $L < L_0$ or $T < T_0$. The bounds are:
\begin{equation}
L_0 \approx \sqrt{\frac{G\hbar}{c^3}} \approx 1.62 \times 10^{-35}\ \text{m}
\end{equation}
\begin{equation}
T_0 \approx \sqrt{\frac{G\hbar}{c^5}} \approx 5.39 \times 10^{-44}\ \text{s}
\end{equation}
These are the \emph{practical} bounds; the absolute bounds (using all energy in the universe) are far smaller but still nonzero.
\end{proof}\TEXTsymbol{\backslash}begin\{proof\}

To measure a length \$L\$, one must use a probe with wavelength \$%
\TEXTsymbol{\backslash}lambda \TEXTsymbol{\backslash}le L\$, requiring
energy \$E \TEXTsymbol{\backslash}ge \TEXTsymbol{\backslash}hbar c/L\$. To
measure a time \$T\$, one requires energy \$E \TEXTsymbol{\backslash}ge 
\TEXTsymbol{\backslash}hbar/T\$. If \$L\$ or \$T\$ is too small, the
required energy exceeds either the total energy of the observable universe
or the threshold for black hole formation. Therefore, no valid measurement
procedure exists for \$L \TEXTsymbol{<} L\_0\$ or \$T \TEXTsymbol{<} T\_0\$.
The bounds are:

\TEXTsymbol{\backslash}begin\{equation\}

L\_0 \TEXTsymbol{\backslash}approx \TEXTsymbol{\backslash}sqrt\{\TEXTsymbol{%
\backslash}frac\{G\TEXTsymbol{\backslash}hbar\}\{c\symbol{94}3\}\} 
\TEXTsymbol{\backslash}approx 1.62 \TEXTsymbol{\backslash}times 10\symbol{94}%
\{-35\}\TEXTsymbol{\backslash} \TEXTsymbol{\backslash}text\{m\}

\TEXTsymbol{\backslash}end\{equation\}

\TEXTsymbol{\backslash}begin\{equation\}

T\_0 \TEXTsymbol{\backslash}approx \TEXTsymbol{\backslash}sqrt\{\TEXTsymbol{%
\backslash}frac\{G\TEXTsymbol{\backslash}hbar\}\{c\symbol{94}5\}\} 
\TEXTsymbol{\backslash}approx 5.39 \TEXTsymbol{\backslash}times 10\symbol{94}%
\{-44\}\TEXTsymbol{\backslash} \TEXTsymbol{\backslash}text\{s\}

\TEXTsymbol{\backslash}end\{equation\}

These are the \TEXTsymbol{\backslash}emph\{practical\} bounds; the absolute
bounds (using all energy in the universe) are far smaller but still nonzero.

\TEXTsymbol{\backslash}end\{proof\}

\subsection{Existence Criterion}

\begin{axiom}[Causal Existence]
An entity $E$ exists in physical reality iff there exists at least one valid
measurement procedure for $E$: 
\begin{equation}
E \in \mathcal{R} \iff \exists\ \text{valid measurement procedure } P \text{
such that } P \text{ detects } E.
\end{equation}
\end{axiom}

\subsection{Practicality}

\begin{axiom}[Finite Resources]
Every interaction in a valid measurement procedure must consume:

\begin{enumerate}
\item Finite energy: $\Delta E_k < \infty$

\item Finite time: $\Delta t_k < \infty$

\item The total energy must be $\le E_{\text{universe}}$, the total energy
of the observable universe
\end{enumerate}
\end{axiom}

\subsection{Wave-Measurement Basis}

\begin{axiom}[Wave-Measurement Basis]
All measurements of length and time are performed using waves (interference
or oscillation). Therefore: 
\begin{equation}
E \ge \max\left(\frac{hc}{\Delta x}, \frac{h}{\Delta t}\right)
\end{equation}
This provides the physical foundation for Axiom 2.
\end{axiom}

\subsection{Information-Reality Separation}

\begin{axiom}[Information-Reality Separation]
Quantum mechanics is a theory of \emph{information about} $\mathcal{R}$, not
of $\mathcal{R}$ itself. The wavefunction, Hilbert space, and operators are
elements of the information layer $\mathcal{I}$, not of the reality layer $%
\mathcal{R}$.
\end{axiom}

\subsection{Dynamic Ontology}

\begin{axiom}[Dynamic Ontology]
The set $\mathcal{R}(t)$ of entities measurable at time $t$ can grow over
time: 
\begin{equation}
t_1 < t_2 \implies \mathcal{R}(t_1) \subseteq \mathcal{R}(t_2)
\end{equation}
However, the criterion for membership (Axioms 1-3) is fixed and invariant.
\end{axiom}

\subsection{Causal Closure}

\begin{axiom}[Causal Closure]
If $E \in \mathcal{R}$, then every physical effect of $E$ is also in $%
\mathcal{R}$: 
\begin{equation}
E \in \mathcal{R} \implies \text{CausalEffect}(E) \subseteq \mathcal{R}.
\end{equation}
\end{axiom}

\subsection{Quantum Logic}

\begin{axiom}[Quantum Logic]
The lattice of measurement propositions is orthomodular and
non-distributive. For propositions $A, B, C$: 
\begin{equation}
A \wedge (B \vee C) \neq (A \wedge B) \vee (A \wedge C).
\end{equation}
This captures the quantum logical structure of yes/no questions about $%
\mathcal{R}$.
\end{axiom}

\section{The Lattice Structure of Measurement Propositions}

\subsection{The Proposition Lattice}

Let $\mathcal{L}$ be the set of all propositions of the form: 
\begin{equation}
\text{"Entity } E \text{ can be measured via procedure } P\text{"}.
\end{equation}

\begin{theorem}
$\mathcal{L}$ forms an orthomodular lattice.
\end{theorem}

\begin{proof}
\begin{proof}
The structure follows from the logical relations among measurement propositions:
\begin{itemize}
    \item \textbf{Meet ($\wedge$)}: Conjunction (AND) of two propositions corresponds to performing both measurements.
    \item \textbf{Join ($\vee$)}: Disjunction (OR) corresponds to performing either measurement.
    \item \textbf{Orthocomplement ($\perp$)}: Negation (NOT) corresponds to the complementary outcome.
\end{itemize}
Orthomodularity follows from the quantum mechanical structure of yes/no questions (Jauch, 1968; Piron, 1976). Non-distributivity arises from incompatible measurements (e.g., position and momentum).
\end{proof}
\end{proof}

\subsection{Non-Distributivity}

The non-distributive nature of $\mathcal{L}$ is captured by: 
\begin{equation}
A \wedge (B \vee C) \neq (A \wedge B) \vee (A \wedge C).
\end{equation}

\begin{example}
Let $A$ be "position in region $X$," $B$ be "momentum in range $P_1$," and $%
C $ be "momentum in range $P_2$." The proposition "position in $X$ and
(momentum in $P_1$ or momentum in $P_2$)" is not equivalent to "(position in 
$X$ and momentum in $P_1$) or (position in $X$ and momentum in $P_2$)"
because the latter requires simultaneous measurement of incompatible
observables.
\end{example}

\subsection{Gleason's Theorem}

\begin{theorem}[Gleason]
Any measure on $\mathcal{L}$ (for $\dim \mathcal{L} \ge 3$) can be
represented by a density matrix $\rho$: 
\begin{equation}
\mu(P) = \func{Tr}(\rho P)
\end{equation}
for all projection operators $P \in \mathcal{L}$.
\end{theorem}

\begin{corollary}
Quantum states (density matrices) are measures on the lattice of measurement
propositions. They are not elements of $\mathcal{R}$; they are elements of
the information layer $\mathcal{I}$.
\end{corollary}

\section{Gravitational Redshift as a Rigorous Test}

The PR's practicality constraint can be tested rigorously using
gravitational redshift. Consider a cesium-133 atom falling toward a black
hole.

\subsection{The Gravitational Redshift Formula}

For a clock at distance $r$ from a spherical mass $M$, an observer at
infinity measures its frequency $f$ as: 
\begin{equation}
f = f_0 \sqrt{1 - \frac{2GM}{rc^2}}
\end{equation}
where $f_0$ is the proper frequency of the clock ($9,192,631,770$ Hz for
cesium-133).

\subsection{The Limit at the Event Horizon}

As $r$ approaches the Schwarzschild radius $r_s = 2GM/c^2$, the term $%
2GM/(rc^2)$ approaches 1, and: 
\begin{equation}
f \rightarrow 0, \qquad E = hf \rightarrow 0
\end{equation}

\begin{theorem}[Event Horizon Limit]
For a cesium atom approaching the event horizon of a black hole, the
observed transition energy $E = hf$ approaches zero, and the measurement of
this energy requires infinite time.
\end{theorem}

\begin{proof}
The gravitational redshift formula gives:
\begin{equation}
E(r) = E_0 \sqrt{1 - \frac{r_s}{r}}
\end{equation}
where $E_0$ is the rest-frame energy. As $r \rightarrow r_s^+$, we have $E(r) \rightarrow 0$. From the perspective of a distant observer, the time between successive ticks of the clock diverges:
\begin{equation}
dt_{\text{obs}} = \frac{dt_{\text{proper}}}{\sqrt{1 - r_s/r}} \rightarrow \infty
\end{equation}
Therefore, measuring the transition energy requires infinite time.
\end{proof}

\subsection{Interpretation Under the PR}

From the perspective of a distant observer:

\begin{itemize}
\item The cesium atom's frequency $f$ approaches zero asymptotically, never
actually reaching zero in finite time.

\item To measure the transition energy going to zero would require infinite
time.

\item Since infinite time is not "practical" or finite, the atomic clock 
\emph{ceases to be a measurable part of our physical reality} as it
approaches the event horizon.
\end{itemize}

This is a direct physical manifestation of the PR: existence is defined by
the possibility of measurement. At the event horizon, the possibility of
measurement breaks down---it would take infinite time---so the entity
effectively ceases to exist in our universe.

\begin{gather}
\text{Gravitational} \\
\begin{tabular}{|c|c|c|}
\hline
\textbf{Parameter} & \textbf{Far from Black Hole} & \textbf{At Event Horizon}
\\ \hline
Observed Frequency $f$ & $\sim 9.19$ GHz & $0$ Hz \\ 
Observed Energy $E=hf$ & Finite & $0$ J \\ 
Measurement Time & Finite (practical) & Infinite (impractical) \\ 
Existence (Per PR) & Exists, measurable & Ceases to exist in measurable
reality \\ \hline
\end{tabular}%
\end{gather}

\section{Quantum Mechanics as Information Theory}

\subsection{The Wavefunction as Knowledge}

In the PR framework, the wavefunction $|\psi\rangle$ represents \emph{%
knowledge about} $\mathcal{R}$, not $\mathcal{R}$ itself. Specifically:

\begin{itemize}
\item $|\psi\rangle$ encodes a probability distribution over possible
measurement outcomes.

\item The dynamics of $|\psi\rangle$ (Schr\"odinger equation) describes the
evolution of knowledge in the absence of measurement.

\item The measurement update (collapse) is an update of knowledge, not a
physical collapse of reality.
\end{itemize}

This is consistent with QBism \cite{Fuchs2010}, Relational Quantum Mechanics 
\cite{Rovelli1996}, and the Brukner-Zeilinger information-theoretic
interpretation \cite{Brukner2001}.

\subsection{The Measurement Problem Resolved}

The measurement problem arises from treating QM as a theory of $\mathcal{R}$%
. If QM is a theory of $\mathcal{I}$, then:

\begin{enumerate}
\item Measurement is \textbf{not} a collapse of reality; it is an update of
information.

\item The wavefunction is a \textbf{preparation procedure}, not an
ontological state.

\item The "collapse" is a \textbf{Bayesian update}: $P(E \mid \text{%
measurement outcome})$ in the lattice $\mathcal{L}$.
\end{enumerate}

Thus, no physical collapse occurs. The measurement problem is dissolved.

\section{Consequences}

\subsection{Planck-Scale Discreteness}

The PR implies that spacetime is discrete at the Planck scale. This is not
an ontological claim about spacetime itself, but a consequence of the
measurability criterion:

\begin{itemize}
\item Lengths below Planck length cannot be measured.

\item Therefore, they are not part of $\mathcal{R}$.

\item The "continuum" of classical spacetime is a mathematical idealization,
not a physical reality.
\end{itemize}

\subsection{The Multiverse}

In the PR framework:

\begin{itemize}
\item Branches of the multiverse that are causally disconnected from our
universe are \textbf{not} in $\mathcal{R}$.

\item They may exist in some abstract sense, but they are \textbf{irrelevant
to physics}.

\item Any branch that \emph{can} affect us is measurable, and therefore
exists in $\mathcal{R}$.
\end{itemize}

\subsection{Hidden Variables}

Hidden variables with no measurable effect are \textbf{not} in $\mathcal{R}$:

\begin{itemize}
\item If a hidden variable has no causal effect on any detector, it cannot
be measured.

\item Therefore, it does not exist in physical reality.

\item It may exist in a formal theory (as a mathematical construct), but it
has no ontological status.
\end{itemize}

\subsection{The "Angels on a Pin" Test}

The PR provides a crisp test for any proposed entity:

\begin{enumerate}
\item Can it be measured via a finite causal chain?

\item Does each step consume finite energy and finite time?

\item If yes to both: it exists in $\mathcal{R}$.

\item If no to either: it does not exist.
\end{enumerate}

This rules out angels, fictional entities, "no intelligent life elsewhere"
claims, and imaginary beings.

\section{Comparison with Other Views}

\begin{table}[h]
\caption{Comparison of Ontological Views}\centering
\begin{tabular}{|c|c|c|c|}
\hline
\textbf{View} & \textbf{What Exists} & \textbf{Role of QM} & \textbf{R/I
Distinction} \\ \hline
Naive Realism & Independent reality & Description of R & No \\ 
Copenhagen & Measurement outcomes & Description of measurement & No \\ 
Relational QM & Relative facts & Description of relative facts & No \\ 
QBism & Agent beliefs & Description of beliefs & I = QM \\ 
\textbf{PR} & \textbf{Causally measurable entities} & \textbf{Description of
I about R} & \textbf{Yes} \\ \hline
\end{tabular}%
\end{table}
\begin{equation}
\begin{tabular}{|c|c|c|c|}
\hline
\textbf{View} & \textbf{What Exists} & \textbf{Role of QM} & \textbf{R/I
Distinction} \\ \hline
Naive Realism & Independent reality & Description of R & No \\ 
Copenhagen & Measurement outcomes & Description of measurement & No \\ 
Relational QM & Relative facts & Description of relative facts & No \\ 
QBism & Agent beliefs & Description of beliefs & I = QM \\ 
\textbf{PR} & \textbf{Causally measurable entities} & \textbf{Description of
I about R} & \textbf{Yes} \\ \hline
\end{tabular}%
\end{equation}

\subsection{Key Differences}

\begin{itemize}
\item \textbf{vs. Naive Realism}: PR does not assume an unobservable reality
beyond measurement.

\item \textbf{vs. Copenhagen}: PR is ontological (existence is defined), not
merely epistemological.

\item \textbf{vs. Relational QM}: PR is non-relative (existence is defined
in our universe).

\item \textbf{vs. QBism}: PR maintains a realist ontology (R exists
independently of agents).
\end{itemize}

\section{Objections and Responses}

\subsection{Objection 1: "The Moon exists when unobserved"}

\textbf{Response:} The Moon is measurable via numerous procedures
(reflecting light, gravitational effects, etc.). Therefore, it is in $%
\mathcal{R}$, even when unobserved. Existence is defined by \emph{%
measurability}, not by current observation.

\subsection{Objection 2: "This is just operationalism"}

\textbf{Response:} Operationalism is about \emph{concepts} (semantic). PR is
about \emph{existence} (ontological). Operationalism: "Length is what we
measure with a ruler." PR: "What we measure with a ruler exists; what cannot
be measured does not."

\subsection{Objection 3: "This leads to solipsism"}

\textbf{Response:} PR defines existence by \emph{possible} measurement, not
actual measurement. The Moon is measurable, so it exists. Solipsism is
avoided.

\subsection{Objection 4: "What about abstract entities (numbers, sets)?"}

\textbf{Response:} Abstract entities are not part of $\mathcal{R}$. They are
part of the mathematical formalism $\mathcal{I}$. They exist in a formal
sense, but they are not physical entities.

\subsection{Objection 5: "What about future discoveries?"}

\textbf{Response:} The PR accommodates future discoveries: new entities can
be added to $\mathcal{R}(t)$ as new measurement procedures are invented. The
criterion for existence remains fixed.

\subsection{Objection 6: "What about entities in disjoint universes?"}

\textbf{Response:} Entities in disjoint universes are causally disconnected
from ours. Therefore, they are not in $\mathcal{R}$. They may exist in some
abstract sense, but they are irrelevant to physics in our universe.

\subsection{Objection 7: "Doesn't this make reality observer-dependent?"}

\textbf{Response:} No. Existence is defined by \emph{possible} measurement
by \emph{any} finite observer, not by actual measurement by a specific
observer. This is objective, not subjective.

\section{Formal Summary}

\subsection{The Principle of Reality (Complete Statement)}

\begin{quote}
\emph{An entity E exists in physical reality iff there exists a finite
sequence of causal interactions:} 
\begin{equation}
E \rightarrow I_1 \rightarrow I_2 \rightarrow \cdots \rightarrow I_n
\rightarrow D
\end{equation}
\emph{where:}
\end{quote}

\begin{itemize}
\item \emph{Each $I_k$ is a physical interaction with finite energy $\Delta
E_k$ and finite duration $\Delta t_k$}

\item \emph{$D$ is a measuring apparatus with a binary (yes/no) output}

\item \emph{$n$ is finite}

\item \emph{The total energy $\sum_k \Delta E_k \le E_{\text{universe}}$}
\end{itemize}

\begin{quote}
\emph{Equivalently: Physical reality is the causal closure of the set of
measurable events.}
\end{quote}

\subsection{The Wave-Measurement Basis}

\begin{equation}
\frame{\boxed{E \ge \max\left(\frac{hc}{\Delta x}, \frac{h}{\Delta t}\right)}%
}
\end{equation}

\subsection{The Information-Reality Distinction}

\begin{equation}
\boxed{ \begin{array}{c} \text{Quantum Mechanics} \\ \text{(Theory of }
\mathcal{I} \text{)} \end{array} \longrightarrow \begin{array}{c}
\text{Physical Reality} \\ (\mathcal{R}) \end{array} }
\end{equation}
\begin{equation}
\text{Wavefunction } |\psi\rangle \in \mathcal{I}, \qquad \text{Entity } E
\in \mathcal{R}
\end{equation}

\subsection{The Dynamic Ontology}

\begin{equation}
\mathcal{R}(t_1) \subseteq \mathcal{R}(t_2) \quad \text{for } t_1 < t_2
\end{equation}
\begin{equation}
\text{Criterion for membership: fixed and invariant}
\end{equation}

\subsection{The Measurement Bound}

\begin{equation}
L_0 \approx \sqrt{\frac{G\hbar}{c^3}} \approx 1.62 \times 10^{-35}\ \text{m}
\end{equation}
\begin{equation}
T_0 \approx \sqrt{\frac{G\hbar}{c^5}} \approx 5.39 \times 10^{-44}\ \text{s}
\end{equation}

\subsection{Causal Closure}

\begin{equation}
E \in \mathcal{R} \implies \text{CausalEffect}(E) \subseteq \mathcal{R}
\end{equation}

\subsection{Non-Boolean Lattice}

\begin{equation}
\mathcal{L} \text{ is orthomodular: } A \wedge (B \vee C) \neq (A \wedge B)
\vee (A \wedge C)
\end{equation}

\subsection{Gravitational Redshift Limit}

\begin{equation}
\lim_{r \to r_s^+} f(r) = 0, \qquad \lim_{r \to r_s^+} E(r) = 0
\end{equation}

\section{Conclusion}

The Principle of Reality provides a consistent, elegant, and physically
grounded ontology for physics:

\begin{enumerate}
\item \textbf{Ontology}: $\mathcal{R}$ is the set of entities connected to
detectors by finite causal chains with finite energy and time.

\item \textbf{Epistemology}: Quantum mechanics is a theory of information
about $\mathcal{R}$, not of $\mathcal{R}$ itself.

\item \textbf{Dynamics}: $\mathcal{R}(t)$ grows over time, but the criterion
for membership is fixed.

\item \textbf{Consequences}: Planck-scale discreteness, causal closure,
non-Boolean lattice structure, resolution of the measurement problem.

\item \textbf{Wave-Measurement Basis}: All measurements are wave
measurements, yielding $E \ge \max(hc/\Delta x, h/\Delta t)$.

\item \textbf{Gravitational Test}: Gravitational redshift provides a
rigorous physical test: transition energies go to zero at the event horizon,
requiring infinite time to measure.
\end{enumerate}

This framework unifies ontology, epistemology, and causality into a single
coherent system. It avoids the pitfalls of both naive realism (which posits
unobservable entities) and antirealism (which denies objective reality). It
provides a clear criterion for what counts as physical reality, and it
accommodates the dynamic nature of scientific discovery.

\begin{acknowledgments}
The author wishes to thank the originator of the Principle of Reality for their deep insights and for provoking this formalization. This work is dedicated to the pursuit of clarity in physics and philosophy.
\end{acknowledgments}

\begin{thebibliography}{9}
\bibitem{Bombelli1987} Bombelli, L., Lee, J., Meyer, D., \& Sorkin, R. D.
(1987). Space-time as a causal set. \textit{Physical Review Letters}, 59(5),
521.

\bibitem{Brukner2001} Brukner, \v{C}., \& Zeilinger, A. (2001). Conceptual
inadequacy of the Shannon information in quantum measurements. \textit{%
Physical Review A}, 63(2), 022113.

\bibitem{Fuchs2010} Fuchs, C. A. (2010). QBism, the perimeter of quantum
Bayesianism. \textit{arXiv:1003.5209}.

\bibitem{Jauch1968} Jauch, J. M. (1968). \textit{Foundations of Quantum
Mechanics}. Addison-Wesley.

\bibitem{Piron1976} Piron, C. (1976). \textit{Foundations of Quantum Physics}%
. W. A. Benjamin.

\bibitem{Rovelli1996} Rovelli, C. (1996). Relational quantum mechanics. 
\textit{International Journal of Theoretical Physics}, 35(8), 1637-1678.

\bibitem{Rovelli1995} Rovelli, C., \& Smolin, L. (1995). Discreteness of
area and volume in quantum gravity. \textit{Nuclear Physics B}, 442(3),
593-619.

\bibitem{Wheeler1990} Wheeler, J. A. (1990). Information, physics, quantum:
The search for links. In \textit{Complexity, Entropy, and the Physics of
Information}. Addison-Wesley.

\bibitem{Zurek2003} Zurek, W. H. (2003). Decoherence, einselection, and the
quantum origins of the classical. \textit{Reviews of Modern Physics}, 75(3),
715.
\end{thebibliography}

\end{document}
